% together/applyrcu.tex
% mainfile: ../perfbook.tex
% SPDX-License-Identifier: CC-BY-SA-3.0

\section{RCU 救援}
\label{sec:together:RCU Rescues}
%
\epigraph{大疑则大悟，小疑则小悟。}
	 {中国谚语}

本节展示如何将 RCU 应用于本书前面讨论的一些示例。
在某些情况下，RCU 提供更简洁的代码；
在另一些情况下提供更好的性能和可扩展性；
还有些情况下两者兼得。

\subsection{RCU 与基于每线程变量的统计计数器}
\label{sec:together:RCU and Per-Thread-Variable-Based Statistical Counters}

\Cref{sec:count:Per-Thread-Variable-Based Implementation}
描述了一种统计计数器的实现，它提供了出色的
性能，大致等同于简单递增（如 C 语言的 \co{++}
运算符），以及线性可扩展性——但仅限于通过
\co{inc_count()} 进行递增。
不幸的是，需要通过 \co{read_count()} 读取值的线程
必须获取全局锁，因此产生了高开销且可扩展性差。
基于锁的实现代码如
\cref{lst:count:Per-Thread Statistical Counters}
（\cpageref{lst:count:Per-Thread Statistical Counters}）所示。

\QuickQuiz{
	我们到底为什么首先需要那个全局锁？
}\QuickQuizAnswer{
	给定线程的 \co{__thread} 变量在该线程退出时消失。
	因此，需要将访问其他线程 \co{__thread} 变量的任何操作
	与线程退出同步。
	如果没有这种同步，访问刚退出线程的 \co{__thread} 变量
	将导致段错误。
}\QuickQuizEnd

\subsubsection{设计}

希望使用 RCU 而非 \co{final_mutex} 来保护
\co{read_count()} 中的线程遍历，
以便从 \co{read_count()} 获得出色的
性能和可扩展性，而不仅仅是从 \co{inc_count()}。
然而，我们不想在计算的总和中牺牲任何精度。
特别是，当给定线程退出时，
我们绝对不能丢失退出线程的计数，
也不能重复计算它。
这样的错误可能导致与结果的全精度相等的不准确性，
换句话说，这样的错误将使结果完全无用。
事实上，\co{final_mutex} 的目的之一就是
确保线程不会在 \co{read_count()} 执行过程中进出。

因此，如果我们要省去 \co{final_mutex}，
就需要想出其他方法来确保一致性。
一种方法是将所有先前退出线程的总计数
和指向每线程计数器的指针数组放入单个结构中。
这样的结构一旦对 \co{read_count()} 可用，
就保持不变，确保 \co{read_count()} 看到一致的数据。

\subsubsection{实现}

\begin{fcvref}[ln:count:count_end_rcu:whole]
\cref{lst:together:RCU and Per-Thread Statistical Counters}
的 \clnrefrange{struct:b}{struct:e}
展示了 \co{countarray} 结构，它包含一个
\co{->total} 字段用于先前退出线程的计数，
以及一个 \co{counterp[]} 指针数组指向每个当前运行线程的
每线程 \co{counter}。
此结构允许 \co{read_count()} 的给定执行
看到与指示的运行线程集一致的总计。

\begin{listing}
\ebresizeverb{.71}{\input{CodeSamples/count/count_end_rcu@whole.fcv}}
\caption{RCU 与每线程统计计数器}
\label{lst:together:RCU and Per-Thread Statistical Counters}
\end{listing}

\Clnrefrange{perthread:b}{perthread:e}
包含每线程 \co{counter} 变量的定义、
引用当前 \co{countarray} 结构的全局指针 \co{countarrayp}，
以及 \co{final_mutex} 自旋锁。

\Clnrefrange{inc:b}{inc:e} 展示了 \co{inc_count()}，
它与 \cref{lst:count:Per-Thread Statistical Counters} 中的相同。
\end{fcvref}

\begin{fcvref}[ln:count:count_end_rcu:whole:read]
\Clnrefrange{b}{e} 展示了 \co{read_count()}，它发生了显著变化。
\Clnref{rrl,rru} 用 \co{rcu_read_lock()} 和
\co{rcu_read_unlock()} 替代了 \co{final_mutex} 的获取和释放。
\Clnref{deref} 使用 \co{rcu_dereference()} 将当前
\co{countarray} 结构快照到局部变量 \co{cap} 中。
正确使用 RCU 将保证此 \co{countarray} 结构
至少在 \clnref{rru} 处当前 RCU 读侧临界区结束之前
一直伴随我们。
\Clnref{init} 将 \co{sum} 初始化为 \co{cap->total}，
即先前退出线程的计数总和。
\Clnrefrange{add:b}{add:e} 累加当前运行线程对应的每线程计数器，
最后，\clnref{ret} 返回总和。
\end{fcvref}

\begin{fcvref}[ln:count:count_end_rcu:whole:init]
\co{countarrayp} 的初始值由
\clnrefrange{b}{e} 的 \co{count_init()} 提供。
此函数在第一个线程创建之前运行，
其任务是分配并清零初始结构，
然后将其赋值给 \co{countarrayp}。
\end{fcvref}

\begin{fcvref}[ln:count:count_end_rcu:whole:reg]
\Clnrefrange{b}{e} 展示了 \co{count_register_thread()} 函数，
它由每个新创建的线程调用。
\Clnref{idx} 获取当前线程的索引，
\clnref{acq} 获取 \co{final_mutex}，
\clnref{set} 安装指向此线程 \co{counter} 的指针，
\clnref{rel} 释放 \co{final_mutex}。
\end{fcvref}

\QuickQuiz{
	\begin{fcvref}[ln:count:count_end_rcu:whole:reg]
	嘿！！！
	\cref{lst:together:RCU and Per-Thread Statistical Counters}
	的 \clnref{set} 修改了一个已有 \co{countarray} 结构中的值！
	你不是说这个结构一旦对 \co{read_count()} 可用就保持不变吗？？？
	\end{fcvref}
}\QuickQuizAnswer{
	我确实是这样说的。
	而且让 \co{count_register_thread()} 分配一个新结构
	是可能的，就像 \co{count_unregister_thread()} 当前所做的那样。

	但这是不必要的。
	回想一下基于内存快照的 \co{read_count()} 误差界限推导。
	因为新线程以初始 \co{counter} 值零开始，
	即使我们在 \co{read_count()} 执行过程中
	添加新线程，推导仍然成立。
	因此，有趣的是，在添加新线程时，
	此实现获得了分配新结构的效果，
	但实际上不必进行分配。
}\QuickQuizEnd

\begin{fcvref}[ln:count:count_end_rcu:whole:unreg]
\Clnrefrange{b}{e} 展示了 \co{count_unregister_thread()}，
它由每个线程在退出前调用。
\Clnrefrange{alloc:b}{alloc:e} 分配一个新的 \co{countarray} 结构，
\clnref{acq} 获取 \co{final_mutex}，\clnref{rel} 释放它。
\Clnref{copy} 将当前 \co{countarray} 的内容复制到新分配的版本中，
\clnref{add} 将退出线程的 \co{counter} 添加到新结构的 \co{->total}，
\clnref{null} 将退出线程的 \co{counterp[]} 数组元素置为 \co{NULL}。
\Clnref{retain} 保留指向当前（即将变为旧的）\co{countarray} 结构的指针，
\clnref{assign} 使用 \co{rcu_assign_pointer()} 安装新版本的 \co{countarray} 结构。
\Clnref{sync} 等待一个宽限期过去，
以便任何可能正在 \co{read_count()} 中并发执行、
从而可能持有对旧 \co{countarray} 结构引用的线程
能够退出其 RCU 读侧临界区，从而释放此类引用。
\Clnref{free} 然后可以安全地释放旧的 \co{countarray} 结构。
\end{fcvref}

\QuickQuiz{
	鉴于固定大小的 \co{counterp} 数组，
	这段代码究竟如何避免对线程数量的固定上限？？？
}\QuickQuizAnswer{
	你说得完全正确，该数组确实重新施加了固定上限。
	可以通过使用链表来跟踪线程来避免此限制，
	就像用户空间 RCU 中所做的那样~\cite{MathieuDesnoyers2012URCU}。
	为此代码做类似的事情留给读者作为练习。
}\QuickQuizEnd

\subsubsection{讨论}

\EQuickQuiz{
	哇！
	\Cref{lst:together:RCU and Per-Thread Statistical Counters}
	包含 70 行代码，而
	\cref{lst:count:Per-Thread Statistical Counters}
	只有 42 行。
	这种额外的复杂性真的值得吗？
}\EQuickQuizAnswer{
	这当然需要根据具体情况来决定。
	如果你需要一个线性可扩展的 \co{read_count()} 实现，
	那么 \cref{lst:count:Per-Thread Statistical Counters}
	中展示的基于锁的实现根本不适合你。
	另一方面，如果对 \co{read_count()} 的调用足够少，
	那么基于锁的版本更简单，因此可能更好，
	尽管大部分大小差异来自结构定义、
	内存分配和 \co{NULL} 返回检查。

	当然，更好的问题是``为什么语言不实现
	对 \co{__thread} 变量的跨线程访问？''
	毕竟，这样的实现将使锁和 RCU 的使用都不必要。
	这将反过来实现一个比
	\cref{lst:count:Per-Thread Statistical Counters}
	中展示的更简单的实现，
	同时具有
	\cref{lst:together:RCU and Per-Thread Statistical Counters}
	中展示的实现的所有可扩展性和性能优势！
}\EQuickQuizEnd

使用 RCU 使退出的线程能够等待，
直到其他线程保证不再使用退出线程的 \co{__thread} 变量。
这允许 \co{read_count()} 函数省去锁定，
从而为 \co{inc_count()} 和 \co{read_count()} 函数
提供出色的性能和可扩展性。
然而，这种性能和可扩展性的代价是代码复杂性的增加。
希望编译器和库的编写者使用用户级
RCU~\cite{MathieuDesnoyers2009URCU}
来提供对 \co{__thread} 变量的安全跨线程访问，
大大降低 \co{__thread} 变量用户所看到的复杂性。

\subsection{RCU 与可移除 I/O 设备的计数器}
\label{sec:together:RCU and Counters for Removable I/O Devices}

\Cref{sec:count:Applying Exact Limit Counters}
展示了一对用于处理可移除设备 I/O 访问计数的
虚构代码片段。
这些代码片段由于需要获取读写锁而在快速路径
（启动 I/O）上产生了高开销。

本节展示如何使用 RCU 来避免此开销。

执行 I/O 的代码与原始版本非常相似，
用 RCU 读侧临界区替代了原始版本中的
读写锁读侧临界区：

\begin{VerbatimN}[tabsize=8]
rcu_read_lock();
if (removing) {
	rcu_read_unlock();
	cancel_io();
} else {
	add_count(1);
	rcu_read_unlock();
	do_io();
	sub_count(1);
}
\end{VerbatimN}
\vspace{5pt}

RCU 读侧原语具有最小开销，因此如期加快了快速路径。

移除设备的更新代码片段如下：

\begin{fcvlabel}[ln:together:applyrcu:Removing Device]
\begin{VerbatimN}[tabsize=8,commandchars=\\\[\]]
spin_lock(&mylock);
removing = 1;
sub_count(mybias);
spin_unlock(&mylock);
synchronize_rcu();
while (read_count() != 0) {	\lnlbl[nextofsync]
	poll(NULL, 0, 1);
}
remove_device();
\end{VerbatimN}
\end{fcvlabel}

\begin{fcvref}[ln:together:applyrcu:Removing Device]
这里我们用独占自旋锁替代读写锁，
并添加 \co{synchronize_rcu()} 来等待所有 RCU 读侧临界区完成。
由于 \co{synchronize_rcu()}，
一旦到达 \clnref{nextofsync}，
我们就知道所有剩余的 I/O 已被计入。

当然，\co{synchronize_rcu()} 的开销可能很大，
但鉴于设备移除相当罕见，这通常是一个好的权衡。
\end{fcvref}

\FloatBarrier
\subsection{数组和长度}
\label{sec:together:Array and Length}

假设我们有一个 RCU 保护的可变长度数组，
如 \cref{lst:together:RCU-Protected Variable-Length Array} 所示。
数组 \co{->a[]} 的长度可以动态变化，
在任何给定时刻，其长度由字段 \co{->length} 给出。
当然，这引入了以下\IX{race condition}：

\begin{listing}
\begin{VerbatimL}[tabsize=8]
struct foo {
	int length;
	char *a;
};
\end{VerbatimL}
\caption{RCU 保护的可变长度数组}
\label{lst:together:RCU-Protected Variable-Length Array}
\end{listing}

\begin{enumerate}
\item	数组最初为 16 个字符长，因此 \co{->length} 等于 16。
\item	CPU~0 加载 \co{->length} 的值，获得值 16。
\item	CPU~1 将数组缩小到长度 8，
	并将指向新的 8 字符内存块的指针赋给 \co{->a[]}。
\item	CPU~0 从 \co{->a[]} 获取新指针，
	并在第 12 个元素中存储新值。
	因为数组只有 8 个字符，这导致
	SEGV 或（更糟的是）内存损坏。
\end{enumerate}

我们如何防止这种情况？

一种方法是小心使用内存屏障，
这在 \cref{chp:Advanced Synchronization: Memory Ordering} 中介绍。
这可以工作，但会产生读侧开销，
而且也许更糟的是需要使用显式内存屏障。

\begin{listing}
\begin{VerbatimL}[tabsize=8]
struct foo_a {
	int length;
	char a[0];
};

struct foo {
	struct foo_a *fa;
};
\end{VerbatimL}
\caption{改进的 RCU 保护的可变长度数组}
\label{lst:together:Improved RCU-Protected Variable-Length Array}
\end{listing}

更好的方法是将值和数组放入同一结构中，
如 \cref{lst:together:Improved RCU-Protected Variable-Length Array}~\cite{Arcangeli03}
所示。
分配新数组（\co{foo_a} 结构）就自动为数组长度提供了新的位置。
这意味着如果任何 CPU 获取了 \co{->fa} 的引用，
\co{->length} 保证与 \co{->a[]} 匹配。

\begin{enumerate}
\item	数组最初为 16 个字符长，因此 \co{->length} 等于 16。
\item	CPU~0 加载 \co{->fa} 的值，
	获得指向包含值 16 和 16 字节数组的结构的指针。
\item	CPU~0 加载 \co{->fa->length} 的值，获得值 16。
\item	CPU~1 将数组缩小到长度 8，
	并将指向包含 8 字符内存块的新 \co{foo_a} 结构的指针
	赋给 \co{->fa}。
\item	CPU~0 从 \co{->a[]} 获取新指针，
	并在第 12 个元素中存储新值。
	但因为 CPU~0 仍在引用包含 16 字节数组的
	旧 \co{foo_a} 结构，一切正常。
\end{enumerate}

当然，在两种情况下，CPU~1 都必须在释放旧数组之前
等待一个宽限期。

下一节将介绍此方法的更通用版本。

\subsection{关联字段}
\label{sec:together:Correlated Fields}
\OriginallyPublished{sec:together:Correlated Fields}{Correlated Fields}{Oregon Graduate Institute}{PaulEdwardMcKenneyPhD}

假设 Sch\"odinger 的每只动物用
\cref{lst:together:Uncorrelated Measurement Fields}
中所示的数据元素表示。
\co{meas_1}、\co{meas_2} 和 \co{meas_3} 字段是
定期更新的一组关联测量值。
读者必须看到来自单次测量更新的这三个值，
这一点至关重要：
如果读者看到 \co{meas_1} 的旧值但 \co{meas_2} 和
\co{meas_3} 的新值，该读者将陷入致命的混乱。
我们如何保证读者将看到这三个值的协调集合？\footnote{
	这种情况类似于
	\cref{sec:together:Correlated Data Elements}
	中描述的情况，
	不同之处在于这里读者只需要看到给定单个数据元素的
	一致视图，而不是前面
	\lcnamecref{sec:together:Correlated Data Elements}
	中所需的一组数据元素的一致视图。}

\begin{listing}
\begin{VerbatimL}[tabsize=8]
struct animal {
	char name[40];
	double age;
	double meas_1;
	double meas_2;
	double meas_3;
	char photo[0]; /* large bitmap. */
};
\end{VerbatimL}
\caption{非关联的测量字段}
\label{lst:together:Uncorrelated Measurement Fields}
\end{listing}

一种方法是分配一个新的 \co{animal} 结构，
将旧结构复制到新结构中，
更新新结构的 \co{meas_1}、\co{meas_2} 和 \co{meas_3} 字段，
然后通过更新指针用新结构替换旧结构。
这确实保证所有读者看到协调的测量值集合，
但由于 \co{->photo[]} 字段，需要复制大型结构。
这种复制可能产生不可接受的高开销。

\begin{listing}
\begin{VerbatimL}[tabsize=8]
struct measurement {
	double meas_1;
	double meas_2;
	double meas_3;
};

struct animal {
	char name[40];
	double age;
	struct measurement *mp;
	char photo[0]; /* large bitmap. */
};
\end{VerbatimL}
\caption{关联的测量字段}
\label{lst:together:Correlated Measurement Fields}
\end{listing}

另一种方法是施加一层间接引用，
如 \cref{lst:together:Correlated Measurement Fields}~\cite[Section 5.3.4]{PaulEdwardMcKenneyPhD}
所示。
当进行新的测量时，分配一个新的 \co{measurement} 结构，
填入测量值，并使用 \co{rcu_assign_pointer()}
更新 \co{animal} 结构的 \co{->mp} 字段指向这个新的
\co{measurement} 结构。
在宽限期过去后，旧的 \co{measurement} 结构可以被释放。

\QuickQuiz{
	但 \cref{lst:together:Correlated Measurement Fields}
	中展示的方法不会导致额外的缓存未命中，
	从而产生额外的读侧开销吗？
}\QuickQuizAnswer{
	确实可能。

\begin{listing}
\begin{VerbatimL}[tabsize=8]
struct measurement {
	double meas_1;
	double meas_2;
	double meas_3;
};

struct animal {
	char name[40];
	double age;
	struct measurement *mp;
        struct measurement meas;
	char photo[0]; /* large bitmap. */
};
\end{VerbatimL}
\caption{本地化的关联测量字段}
\label{lst:together:Localized Correlated Measurement Fields}
\end{listing}

	避免此缓存未命中开销的一种方法如
	\cref{lst:together:Localized Correlated Measurement Fields} 所示：
	只需在 \co{animal} 结构中嵌入一个名为 \co{meas} 的
	\co{measurement} 结构实例，
	并将 \co{->mp} 字段指向此 \co{->meas} 字段。

	然后可以按以下步骤进行测量更新：

	\begin{enumerate}
	\item	分配一个新的 \co{measurement} 结构，
		并将新的测量值放入其中。
	\item	使用 \co{rcu_assign_pointer()} 将 \co{->mp}
		指向此新结构。
	\item	等待宽限期过去，例如使用
		\co{synchronize_rcu()} 或 \co{call_rcu()}。
	\item	将新 \co{measurement} 结构中的测量值
		复制到嵌入的 \co{->meas} 字段中。
	\item	使用 \co{rcu_assign_pointer()} 将 \co{->mp}
		指回旧的嵌入 \co{->meas} 字段。
	\item	在另一个宽限期过去后，
		释放新的 \co{measurement} 结构。
	\end{enumerate}

	这种方法使用较重的更新过程来消除
	常见情况下的额外缓存未命中。
	额外的缓存未命中仅在更新实际进行中时才会发生。
}\QuickQuizEnd

这种方法使读者能够看到选定字段的关联值，
同时产生最小的读侧开销。
在读者仅查看单个数据元素的常见情况下，
这种每数据元素的一致性就已足够。

% Flag for deletion (if not already covered in the defer chapter).
% @@@ Issaquah Challenge.
% @@@ RLU & MV-RLU (Eventually the corresponding patents.)

\subsection{更新友好型遍历}
\label{sec:together:Update-Friendly Traversal}

假设需要对 hash table 中所有元素进行统计扫描。
例如，Schr\"odinger 可能希望计算其所有动物的
平均体长体重比。\footnote{
	为什么这样的量会有用？
	我也不知道！
	但分组统计通常是有用的。}
进一步假设 Schr\"odinger 愿意忽略
由于在统计扫描进行期间动物被添加到
hash table 或从中删除而导致的轻微误差。
Schr\"odinger 应该怎样控制并发？

一种方法是将统计扫描封装在 RCU 读侧临界区中。
这允许更新在不过度阻碍扫描的情况下并发进行。
特别是，扫描不会阻塞更新，反之亦然，
这使得即使面对非常高的更新速率，
也能优雅地支持包含大量元素的 hash table 扫描。

\QuickQuiz{
	但在调整 hash table 大小时，此扫描如何工作？
	在那种情况下，旧的和新的 hash table
	都不能保证包含 hash table 中的所有元素！
}\QuickQuizAnswer{
	确实，如
	\cref{sec:datastruct:Non-Partitionable Data Structures}
	中描述的可调整大小的 hash table
	在调整大小期间不能被完全扫描。
	一个简单的解决方法是在扫描时获取
	\co{hashtab} 结构的 \co{->ht_lock}，
	但这会阻止多个扫描同时进行。

	另一种方法是让更新在调整大小进行中时
	同时修改旧的和新的 hash table。
	这将允许扫描在旧 hash table 中找到所有元素。
	实现这一点留给读者作为练习。
}\QuickQuizEnd

\subsection{可扩展引用计数之二}
\label{sec:together:Scalable Reference Count Two}

假设\IX{reference count}正在成为性能或可扩展性的瓶颈。
你能做什么？

一种方法是为每个引用计数使用每 CPU 计数器，
与 \cref{chp:Counting} 中的算法类似，
特别是 \cref{sec:count:Exact Limit Counters}
中描述的精确限制计数器。
需要在这些计数器的每 CPU 和全局模式之间切换，
这导致要么递增和递减开销高
（\cref{sec:count:Atomic Limit Counter Implementation}），
要么使用 POSIX 信号
（\cref{sec:count:Signal-Theft Limit Counter Design}）。

另一种方法是使用 RCU 来调解每 CPU 和全局计数模式之间的切换。
每个更新在 RCU 读侧临界区内进行，
每个更新检查一个标志以确定是更新
每 CPU 计数器还是全局计数器。
要切换模式，更新标志，等待宽限期，
然后将任何剩余的计数从每 CPU 计数器
移到全局计数器，反之亦然。

Linux 内核在其 \co{percpu_ref} 引用计数器风格中使用了
这种 RCU 中介的方法。
使用此引用计数器的代码必须使用 \co{percpu_ref_init()}
初始化 \co{percpu_ref} 结构，
该函数接受的参数包括：指向结构的指针、
引用计数达到零时调用的函数的指针、
一组模式标志和一组 \co{kmalloc()} \co{GFP_} 标志。
正常初始化后，结构具有一个引用并处于每 CPU 模式。

模式标志通常为零，但如果计数器要从
慢速非每 CPU（即原子）模式开始，
可以包含 \co{PERCPU_REF_INIT_ATOMIC} 位。
还有一个 \co{PERCPU_REF_ALLOW_REINIT} 位，
允许给定的 \co{percpu_ref} 计数器通过调用
\co{percpu_ref_reinit()} 来重用，
而无需释放和重新分配。
无论 \co{percpu_ref} 结构如何初始化，
都可以使用 \co{percpu_ref_get()} 获取引用，
使用 \co{percpu_ref_put()} 释放引用。

在每 CPU 模式下，\co{percpu_ref} 结构
无法确定其值是否已达到零。
当需要这样的判断时，可以调用 \co{percpu_ref_kill()}。
此函数将结构切换到原子模式，
并删除由 \co{percpu_ref_init()} 调用安装的初始引用。
当然，在原子模式下，调用 \co{percpu_ref_get()} 和
\co{percpu_ref_put()} 相当昂贵，
但 \co{percpu_ref_put()} 可以知道值何时达到零。

希望了解更多 \co{percpu_ref} 信息的读者
请参阅 Linux 内核文档和源代码。

\subsection{重触发宽限期}
\label{sec:together:Retriggered Grace Periods}

没有 \co{call_rcu_cancel()}，
所以一旦 \co{rcu_head} 结构传递给 \co{call_rcu()}，
就无法收回。
必须让它保持不变，直到回调被调用。
在常见情况下，这是应该的，
因为 \co{rcu_head} 结构正走在通往释放的单行道上。

然而，有些用例将 RCU 与显式的 \co{open()} 和
\co{close()} 调用结合使用。
在 \co{close()} 调用之后，读者不应该开始
对数据结构的新访问，
但很可能有读者正在完成其遍历。
这种情况可以用通常的方式处理：
在 \co{close()} 调用之后等待一个宽限期，
然后再释放数据结构。

但如果在宽限期结束前调用了 \co{open()} 怎么办？

同样，没有 \co{call_rcu_cancel()}，
所以另一种方法是设置一个由回调函数检查的标志，
回调函数可以选择不实际释放任何东西。
问题解决了！

但如果在宽限期结束前先调用了 \co{open()}
然后又调用了 \co{close()} 怎么办？

一种方法是为标志设置第二个值，
使回调重新将自身排队。

但如果不仅有 \co{open()} 然后 \co{close()}，
还有在宽限期结束前的另一个 \co{open()} 呢？

在这种情况下，回调需要设置状态来反映
最后一个 \co{open()} 仍然有效。

\begin{figure}
\centering
\resizebox{2.5in}{!}{\includegraphics{together/retriggergp}}
\caption{重触发宽限期状态机}
\label{fig:count:Retrigger-Grace-Period State Machine}
\end{figure}

沿着这条思路继续，我们得到了
\cref{fig:count:Retrigger-Grace-Period State Machine}
中所示的状态机。
初始状态是 CLOSED，运行状态是 OPEN\@。
菱形箭头表示 \co{call_rcu()} 调用，
标有 ``CB'' 的箭头表示回调调用。

通过此状态机的正常路径遍历状态 CLOSED、
OPEN、CLOSING（调用 \co{call_rcu()}），
然后在回调被调用后回到 CLOSED。
如果在宽限期完成之前调用了 \co{open()}，
状态机遍历循环 OPEN、CLOSING（调用
\co{call_rcu()}）、REOPENING，
然后在回调被调用后回到 OPEN。
如果在宽限期完成之前调用了 \co{open()} 然后 \co{close()}，
状态机遍历循环 OPEN、CLOSING（调用
\co{call_rcu()}）、REOPENING、RECLOSING，
然后在回调被调用后回到 CLOSING。

给定一个无限的 \co{close()} 和 \co{open()} 交替调用序列，
状态机将遍历 OPEN 和 CLOSING（调用 \co{call_rcu()}），
然后交替停留在 REOPENING 和 RECLOSING 状态。
一旦宽限期结束，状态机将转换到
CLOSING 或 OPEN 状态，
取决于回调是在 RECLOSING 还是 REOPENING 状态中被调用。

\begin{listing}
\ebresizeverb{.88}{\input{CodeSamples/together/retrigger-gp@whole.fcv}}
\caption{重触发宽限期（伪代码）}
\label{lst:together:Retriggering a Grace Period}
\end{listing}

此状态机的粗略伪代码如
\cref{lst:together:Retriggering a Grace Period} 所示。
\begin{fcvref}[ln:together:retrigger-gp:whole]
五个状态如 \clnrefrange{states:b}{states:e} 所示，
当前状态保存在 \clnref{status} 的 \co{rtrg_status} 中，
由 \clnref{lock} 的锁保护。

三个 CB 转换（从状态 CLOSING、REOPENING
和 RECLOSING 发出）由
\clnrefrange{close_cb:b}{close_cb:e} 的 \co{close_cb()} 函数实现。
\Clnref{cleanup} 在转换到 CLOSED 状态时
调用用户提供的 \co{close_cleanup()}
来执行任何最终的清理操作，如释放内存。
\Clnref{call_rcu1} 包含 \co{call_rcu()} 调用，
它导致稍后转换到 CLOSED 状态。

\clnrefrange{open:b}{open:e} 的 \co{open()} 函数实现
到 OPEN、CLOSING 和 REOPENING 状态的转换，
\clnref{do_open} 调用 \co{do_open()} 函数来实现
任何所需数据结构的分配和初始化。

\clnrefrange{close:b}{close:e} 的 \co{close()} 函数
实现到 CLOSING 和 RECLOSING 状态的转换，
\clnref{do_close} 调用 \co{do_close()} 函数来执行
完成此转换可能需要的任何操作，
例如导致后续的只读遍历返回错误。
\Clnref{call_rcu2} 包含 \co{call_rcu()} 调用，
它导致稍后转换到 CLOSED 状态。
\end{fcvref}

此状态机和伪代码展示了如何在那些罕见的
需要此类语义的情况下获得 \co{call_rcu_cancel()} 的效果。

\subsection{长时间访问之二}
\label{sec:together:Long-Duration Accesses Two}

假设一个读写锁的读者持有锁的时间太长，
导致更新被过度延迟。
进一步假设该读者无法合理地转换为
使用引用计数
（否则参见 \cref{sec:together:Long-Duration Accesses}）。

如果该读者可以合理地转换为使用 RCU，
这可能解决问题。
原因是 RCU 读者不会完全阻塞更新，
而只阻塞那些更新的清理部分（包括内存回收）。
因此，如果系统有充足的内存，
将读写锁转换为 RCU 可能就足够了。

然而，转换为 RCU 并不总是足够的。
例如，代码可能在单个 RCU 读侧临界区内
遍历一个极其大的链式数据结构，
这可能极大地延长 RCU 宽限期，导致系统内存耗尽。
这些情况可以用几种不同的方式处理：
\begin{enumerate*}[(1)]
\item	使用 SRCU 代替 RCU，以及
\item	获取引用以退出 RCU 读者。
\end{enumerate*}

\subsubsection{使用 SRCU}
\label{sec:together:Use SRCU}

在 Linux 内核中，RCU 是全局的。
换句话说，内核中任何地方的长时间运行的 RCU 读者
都会延迟当前的 RCU 宽限期。
如果长时间运行的 RCU 读者正在遍历一个小数据结构，
那少量的数据就延迟了所有其他数据结构的释放，
这可能导致内存耗尽。

避免此问题的一种方法是对该长时间运行的
RCU 读者的数据结构使用 SRCU，
使用其自己的 \co{srcu_struct} 结构。
由此产生的长时间运行的 SRCU 读者
将仅延迟该 \co{srcu_struct} 结构的宽限期，
而不延迟 RCU 的宽限期，
从而避免内存耗尽。
更多详情请参见 \cref{sec:defer:RCU Linux-Kernel API} 中的 SRCU API。

不幸的是，这种方法确实有一些缺点。
首先，SRCU 读者不受优先级提升的约束，
这可能导致繁忙系统上低优先级 SRCU 读者的额外延迟。
更糟的是，定义单独的 \co{srcu_struct} 结构
减少了 RCU 更新者的数量，
这反过来增加了每个更新者的宽限期开销。
这意味着为每个当前 Linux 内核 RCU 用例
都提供其自己的 \co{srcu_struct} 结构，
可能会将系统范围的宽限期开销
乘以此类结构的数量。

因此，通常更好的做法是获取某种非 RCU 引用
以暂时退出 RCU 读侧临界区，
如下一节所述。

\subsubsection{获取引用}
\label{sec:together:Acquire a Reference}

如果 RCU 读侧临界区太长，就缩短它！

在某些情况下，这可以轻松完成。
例如，扫描静态分配的哈希桶数组的所有哈希链的代码
可以同样轻松地在各自的临界区中扫描每个哈希链。

这之所以有效，是因为哈希链通常相当短，而且是设计如此。
当遍历长链式结构时，必须有某种方式
在中间停下来并在稍后恢复。

例如，在 Linux 内核 v5.16 中，
\co{khugepaged_scan_file()} 函数使用 \co{need_resched()}
检查其他任务是否需要当前 CPU，
如果是，则调用 \co{xas_pause()} 适当调整
遍历的迭代器，然后调用 \co{cond_resched_rcu()}
让出 CPU\@。
\co{cond_resched_rcu()} 函数依次调用 \co{rcu_read_unlock()}、
\co{cond_resched()} 和最后 \co{rcu_read_lock()}，
以退出 RCU 读侧临界区来让出 CPU。

当这种简单的临界区缩短不可行时，
另一种方法是获取对象的引用，
然后退出 RCU 读侧临界区。
v6.17 内核中的 \co{txopt_get()} 函数采用了这种方法，
使用 \co{refcount_inc_not_zero()} 尝试获取引用。
如果由此产生的内存竞争过大，
另一个修复方法是使用 hazard pointer 而非显式引用计数。

当然，在可行的情况下，另一种方法是切换到
对暂时退出 RCU 读侧临界区更友好的
数据结构，如 hash table。

\QuickQuiz{
	但这如何与可调整大小的 hash table 一起工作，
	例如
	\cref{sec:datastruct:Non-Partitionable Data Structures}
	中描述的那个？
}\QuickQuizAnswer{
	在这种情况下，需要更多的注意，
	因为在我们暂时退出 RCU 读侧临界区期间，
	hash table 很可能被调整了大小。
	更糟的是，调整大小操作可以预期会释放旧的哈希桶，
	使我们指向空闲列表。

	但仅仅阻止旧哈希桶被释放是不够的。
	还必须确保那些桶继续被更新。

	处理此问题的一种方法是在每组桶上
	设置一个引用计数，初始值设为 1。
	全表扫描将在扫描开始时获取一个引用
	（但仅在引用非零时），
	并在扫描结束时释放它。
	调整大小操作将填充新桶、释放引用、
	等待宽限期，然后等待引用变为零。
	一旦引用为零，调整大小操作可以让更新者
	忘记旧的哈希桶，然后释放它。

	实际实现留给有兴趣的读者，
	他们将从此任务中获得很多见解。
}\QuickQuizEnd

\subsection{无锁配置更新}
\label{sec:together:Lockless Configuration Update}

假设你有一个全局指针，
它要么是 \co{NULL}，
要么指向一个提供配置信息的单例数据结构，
类似于
\cref{sec:defer:Minimal Insertion and Deletion} 中讨论的情况。
本节增加的变化是允许并发无锁更新。

\begin{listing}
\input{CodeSamples/formal/herd/C-MP+o-xchg-rcusync-o+o-xchg-rcusync-o+rl-o-o-rul@whole.fcv}
\caption{配置更新示例}
\label{lst:formal:Configuration-Update Example}
\end{listing}

乍一看，该节中的更新是无锁的，
因为它们使用 \co{rcu_assign_pointer()}。
然而，如果我们有并发的无锁更新，
一个更新的 \co{rcu_assign_pointer()}
可能被下一个更新悄悄覆盖，导致内存泄漏。
\begin{fcvref}[ln:formal:C-MP+o-xchg-rcusync-o+o-xchg-rcusync-o+rl-o-o-rul:whole]
避免这种情况的一种方法是使用原子 \co{xchg()} 指令，
如 \cref{lst:formal:Configuration-Update Example}
的 \clnref{P0:xchg, P1:xchg} 所示。
这样，被覆盖的 \co{xchg()} 返回的指针
由下一个 \co{xchg()} 返回，
允许调用者在 RCU 宽限期后释放它
（\clnref{P0:sync, P1:sync}）。
因为 Linux 内核内存模型不支持动态分配，
此代码通过存储值零来模拟释放。
\end{fcvref}

\QuickQuizSeries{%
\QuickQuizB{
	等等！！！
	\begin{fcvref}[ln:formal:C-MP+o-xchg-rcusync-o+o-xchg-rcusync-o+rl-o-o-rul:whole]
	\cref{lst:formal:Configuration-Update Example}
	的 \clnref{P0:free, P1:free}
	不需要检查从 \clnref{P0:xchg, P1:xchg} 返回的
	\co{NULL} 指针吗？
	\end{fcvref}
}\QuickQuizAnswerB{
	在这个教科书示例中，不需要。

	\begin{fcvref}[ln:formal:C-MP+o-xchg-rcusync-o+o-xchg-rcusync-o+rl-o-o-rul:whole]
	这是因为 \clnref{init:p} 将指针 \co{p} 初始化为
	指向变量 \co{a}，所以此指针永远不会是 \co{NULL}。
	然而，这是一个很好的问题，
	因为在大多数生产级代码中，
	确实需要 \co{NULL} 检查。
	此外，本节的第一句话暗示指针实际上可能是 \co{NULL}。
	添加 NULL 检查留给读者作为练习。
	\end{fcvref}
}\QuickQuizEndB
%
\QuickQuizE{
	我们在 \cref{lst:formal:Configuration-Update Example}
	中真的需要 \co{xchg()} 吗？
	为什么不用简单的加载后跟简单的存储？
}\QuickQuizAnswerE{
	是的，我们确实需要那个 \co{xchg()}。
	如果我们使用简单的加载后跟简单的存储，
	另一个进程可能在加载和存储之间更新指针，
	导致该进程的元素被覆盖和泄漏。

	但不要只听\emph{我}的话。
	用 \co{herd} 试试看！！！
	(\path{C-MP+o-o-rcusync-o+o-o-rcusync-o+rl-o-o-rul.litmus})
}\QuickQuizEndE
}		% End of \QuickQuizSeries

\begin{fcvref}[ln:formal:C-MP+o-xchg-rcusync-o+o-xchg-rcusync-o+rl-o-o-rul:whole]
读者然后可以作为直接的 RCU 读者编写
（\clnrefrange{P2:lock}{P2:unlock}），
在 \clnref{P2:load} 加载当前指针，
在 \clnref{P2:deref} 解引用它。

运行此示例
(\path{CodeSamples/formal/herd/C-MP+o-xchg-rcusync-o+o-xchg-rcusync-o+rl-o-o-rul.litmus})
通过 \co{herd}
（参见 \cref{sec:formal:Axiomatic Approaches}）
验证了此算法避免了释放后使用错误（\clnref{ex:1}），
并且它模拟了释放对象 \co{a}、\co{b} 和 \co{c}
中除一个以外的所有对象（\clnref{ex:2}），
从而也避免了内存泄漏。
\end{fcvref}

\subsection{无锁双重检查锁定}
\label{sec:together:Lockless Double-Checked Locking}

双重检查锁定是一种众所周知的模式，
提供共享状态的首次使用初始化。
但考虑一个极端低延迟环境，
其中禁止任何和所有锁的使用。
在这种约束下如何实现首次使用初始化（和后续更新）？

\begin{listing}
\ebresizeverb{.9}{
\input{CodeSamples/formal/herd/C-double-check-rcu-2@whole.fcv}
}
\caption{两进程双重检查 RCU}
\label{lst:formal:Two-Process Double-Check RCU}
\end{listing}

一种方法是使用 RCU 来保护查找操作，
并使用原子 \co{xchg()} 操作来执行更新。
与上一节一样，被覆盖的对象通过 RCU 释放
以避免干扰并发读者，例如
\cref{lst:formal:Two-Process Double-Check RCU}
(\path{C-double-check-rcu-2.litmus}) 所示。\footnote{
	对两进程示例不信服的人可以查看
	三进程示例
	(\path{C-double-check-rcu-3.litmus})。
	对三进程示例不信服的人可以构建更大的示例，
	但请注意在 Paul 的笔记本电脑上，
	\co{herd} 分析四进程示例大约需要 20 分钟，
	而三进程示例大约需要 2 秒。}

快速路径覆盖了初始化已经完成的常见情况。
\begin{fcvref}[ln:formal:C-double-check-rcu-2:whole:P0]
\Clnref{lock} 开始 RCU 读侧临界区，
\clnref{load} 获取当前指针。
如果 \clnref{if} 确定此指针非 \co{NULL}，
则初始化已完成，
\clnref{deref} 获取元素的内容，
\clnref{unl:1} 退出临界区。

否则，\clnref{if} 将看到初始化前的 \co{NULL} 指针，
并将控制转移到 \clnref{unl:2}，退出 RCU 读侧临界区。
\Clnref{init:a} 初始化元素（抽象掉任何所需的分配），
\clnref{xchg} 将指向此元素的指针与全局指针 \co{p}
原子交换，将旧值返回到 \co{r1} 中。
如果 \clnref{if:2} 确定 \co{r1} 非 \co{NULL}，
则 \clnref{sync} 等待任何访问旧元素的读者完成，
\clnref{free} 模拟 \co{kfree()}。
无论如何，\co{r2} 被设置为初始值 1。
\end{fcvref}

\begin{fcvref}[ln:formal:C-double-check-rcu-2:whole]
\Clnrefrange{P1:b}{P1:e} 以相同方式操作，
但使用元素 \co{b} 而非 \co{a}。

\Clnref{loc} 显示所有变量以便于调试 litmus 测试。
\Clnref{ex:1} 验证两个进程都获得了初始化值，
\clnref{ex:2} 验证没有内存泄漏且至少有一个元素仍在使用中。
运行 \co{herd} 确认此代码满足这些约束。
\end{fcvref}

\QuickQuiz{
	但在这样的低延迟环境中，
	内存分配不也是被禁止的吗？
}\QuickQuizAnswer{
	确实可能。
	在这种情况下，可能需要预分配，
	例如在 Linux 内核中，
	每个 CPU 可能被提供一个预分配的元素。
}\QuickQuizEnd

\subsection{无锁双重检查初始化}
\label{sec:together:Lockless Double-Checked Initialization}

假设我们有
\cref{sec:together:Lockless Double-Checked Locking}
中描述的情况，但不需要后续更新。
在这种约束下实现首次使用初始化（但不需要后续更新）
的最佳方法是什么？

\begin{listing}
\ebresizeverb{.9}{
\input{CodeSamples/formal/herd/C-double-check-cas-2@whole.fcv}
}
\caption{两进程双重检查 CAS}
\label{lst:formal:Two-Process Double-Check CAS}
\end{listing}

在这种情况下，我们只需设置指针 \co{p} 一次。
这意味着如果任何想要初始化的进程发现
此指针非 \co{NULL}，它可以简单地放弃其初始化尝试。
这反过来意味着任何读者只会获得对
唯一（永生）元素的引用，
消除了对任何延迟回收的需要。
因此，可以使用简单的 \co{cmpxchg()} 来实现此初始化过程，
如 \cref{lst:formal:Two-Process Double-Check CAS} 所示。

\begin{fcvref}[ln:formal:C-double-check-cas-2:whole:P0]
\Clnref{load} 获取指针，如果 \clnref{if} 看到它
非 \co{NULL}，\clnref{deref} 获取值。

否则，\clnref{init:a} 初始化一个元素，
\clnref{cmpxchg} 尝试将其安装在指针 \co{p} 处。
如果此尝试失败，意味着其他进程已经完成了初始化，
所以 \clnref{free} 释放此进程的元素。
\end{fcvref}

\QuickQuiz{
	\begin{fcvref}[ln:formal:C-double-check-cas-2:whole:P0]
	\cref{lst:formal:Two-Process Double-Check CAS}
	的 \clnref{load} 上的 \co{READ_ONCE()}
	不需要是 \co{smp_load_acquire()} 以防止
	与初始化并发运行的早期读者
	看到初始化前的垃圾数据吗？
	\end{fcvref}
}\QuickQuizAnswer{
	不需要，因为 \co{READ_ONCE()} 是 volatile 操作，
	因此可以作为地址依赖链的头部。
	理论上，另一个替代方案是 \co{rcu_dereference_raw()}。
	除了这可能会产生误导，
	因为该元素永远不会被释放，
	不受 \co{rcu_read_lock()} 保护，
	因此不是真正的 RCU 保护指针。

	有关地址依赖的更多信息，请参阅
	\cref{sec:memorder:Address- and Data-Dependency Difficulties}
	（\cpageref{sec:memorder:Address- and Data-Dependency Difficulties}
	开始）。
}\QuickQuizEnd

此示例强调了彻底理解需求的必要性。
在此示例中，删除运行时更新的需求
通过消除对 RCU 的需求简化了同步。\footnote{
	感谢 Denis Yaroshevskiy 指出这一点。}

\QuickQuiz{
	给定如
	\cref{lst:formal:Two-Process Double-Check CAS}
	所示的算法，
	为什么还有人使用双重检查锁定？
}\QuickQuizAnswer{
	如果标志和数据被仔细放入同一缓存行，
	那么双重检查锁定可以提供比双重检查 CAS
	更好的缓存局部性。
	这是因为在双重检查 CAS 情况下，
	指针和被指向的元素几乎总是在不同的缓存行中，
	可能导致两次缓存未命中，
	而双重检查锁定只有一次缓存未命中。

	然而，随着元素大小的增长，这个优势会减弱。
}\QuickQuizEnd
